\begingroup
\renewcommand{\arraystretch}{1.45}
\Needspace{37\baselineskip}
\InterpolationArticle{15}
Aus dem vorangehenden Satz folgt: Sind die Werte einer Funktion $X$ von der Form $F$ für $x=a$, $b$, $c$, $d$ usw. der Reihe nach $A$, $B$, $C$, $D$ usw., so muss jede andere ähnliche Funktion des Bogens $x$, die denselben Werten genügt, in der Formel $X+PY$ enthalten sein, wobei wir mit $P$ das Produkt
\[
\sin\tfrac12(x-a)\sin\tfrac12(x-b)\sin\tfrac12(x-c)\sin\tfrac12(x-d)\quad\text{usw.}
\]
bezeichnen, mit $Y$ dagegen eine unbestimmte Funktion des Bogens $x$, die von der Form $F$ oder von der Form
\SourceLeaf{298}
\(G\) ist, je nachdem, ob die Anzahl der Werte \(a,b,c,d\) usw., die wir mit \(\mu\) bezeichnen werden, gerade oder ungerade ist. Wir bemerken nun:

I. In den Artikeln 10, 11 und 12 haben wir gezeigt, wie eine Funktion \(X\) von \(m^{\mathrm{ter}}\) Ordnung bestimmt wird (die also nicht über \(\cos mx\) und \(\sin mx\) hinausgeht), die den \(\mu\) vorgegebenen Werten \(A,B,C,D\) usw. genügt, wobei \(m\) \(=\frac{1}{2}\mu-\frac{1}{2}\) oder \(\frac{1}{2}\mu\) ist, je nachdem, ob \(\mu\) ungerade oder gerade ist. Das Produkt \(P\) ist offenbar von \({\frac{1}{2}\mu}^{\mathrm{ter}}\) Ordnung. Ist daher \(\mu\) ungerade und somit \(Y\) von \({\frac{1}{2}}^{\mathrm{ter}}\) oder höherer Ordnung, so ist \(PY\) mindestens von \({\frac{1}{2}\mu+\frac{1}{2}}^{\mathrm{ter}}\) Ordnung und folglich von höherer Ordnung als \(X\). Daraus folgt, dass andere Funktionen als \(X\), die denselben vorgegebenen Werten genügen, nur von höherer Ordnung sein können. Ist dagegen \(\mu\) gerade, so ist \(PY\) nur dann von \({\frac{1}{2}\mu}^{\mathrm{ter}}\) Ordnung, wenn für \(Y\) eine bestimmte Größe genommen wird; in diesem Fall gibt es also zwar unendlich viele Funktionen von ähnlicher Form und gleicher Ordnung wie \(X\), die denselben Werten genügen, doch sind sie alle in der Form \(X+hP\) enthalten, wobei \(h\) eine bestimmte Größe bezeichnet. Dieselben Folgerungen ergeben sich unmittelbar aus der Methode, mit der wir in den Artikeln 10 und 11 die Funktion \(X\) hergeleitet haben.

\Needspace{26\baselineskip}
II. Ist umgekehrt von anderswoher eine Funktion \(X'\) bekannt, die zwar allen vorgegebenen Werten \(A,B,C,D\) usw. genügt, aber von höherer Ordnung ist als nötig, so lässt sich die Funktion \(Y\) so bestimmen, dass \(X'+PY\) auf \({\frac{1}{2}\mu-\frac{1}{2}}^{\mathrm{te}}\) oder \({\frac{1}{2}\mu}^{\mathrm{te}}\) Ordnung herabgesetzt wird (je nachdem, ob \(\mu\) ungerade oder gerade ist). Die höchsten Glieder in \(X'\) seien \(K\cos nx+L\sin nx\), diejenigen in \(P\) seien \(k\cos\frac{1}{2}\mu x+l\sin\frac{1}{2}\mu x\); als höchste Glieder in \(Y\) nehme man
\[
\rho=-\frac{2(Kk+Ll)}{kk+ll}\cos(n-\tfrac{1}{2}\mu)x
     +\frac{2(Kl-Lk)}{kk+ll}\sin(n-\tfrac{1}{2}\mu)x
\]
Durch Rechnung ergibt sich dann, dass die Koeffizienten von \(\cos nx\) und \(\sin nx\) in \(X'+PY\) \(=0\) werden, sofern \(n>\frac{1}{2}\mu\) gilt; ebenso werden aus den höchsten Gliedern der entwickelten Funktion \(X'+P\rho\) (die von \({n-1}^{\mathrm{ter}}\) Ordnung sein werden) die nächstniedrigeren Glieder in \(Y\) (von \({n-\frac{1}{2}\mu-1}^{\mathrm{ter}}\) Ordnung) bestimmt. Diese Operation ist so lange zu wiederholen, bis \(X'+PY\) auf eine nicht höhere Ordnung als \(\frac{1}{2}\mu\) herabgesetzt ist, also auf \({\frac{1}{2}\mu-\frac{1}{2}}^{\mathrm{te}}\) Ordnung, wenn \(\mu\) ungerade ist, und auf \({\frac{1}{2}\mu}^{\mathrm{te}}\) Ordnung, wenn \(\mu\) gerade ist. Dass sich diese Herabsetzung nicht weiter fortführen lässt, ist ohne Schwierigkeit einzusehen.

\Needspace{46\baselineskip}
III. Im ersten Fall (wenn \(\mu\) ungerade ist) werden wir auf diese Weise also notwendig dieselbe Funktion \(X\) erhalten, die wir nach der oben angegebenen Methode aus den gegebenen Werten \(A,B,C,D\) usw. bestimmt hätten. Im zweiten Fall liefert die eben erläuterte Methode dagegen eine Funktion, die zwar von derselben Ordnung ist wie die Funktion, die nach der Methode von Artikel 12 oder nach der Methode von Artikel 11 entsteht, wenn das letzte Glied
\SourceLeaf{299}
\(=0\) gesetzt wird, die aber von dieser verschieden sein kann. Die erstere Funktion sei \(X''\), die letztere \(X'''\); also gilt \(X'''=X''+hP\), wobei \(h\) eine bestimmte Größe bezeichnet. Wir setzen als höchste Glieder in \(X''\) \(K'\cos\frac{1}{2}\mu x+L'\sin\frac{1}{2}\mu x\), in \(X'''\) dagegen \(K''\cos\frac{1}{2}\mu x+L''\sin\frac{1}{2}\mu x\) an, sodass \(K''=K'+hk\), \(L''=L'+hl\) gilt. Nun lässt sich aus Artikel 11 oder Artikel 12 ohne Schwierigkeit zeigen: Bezeichnet man die Summe \(a+b+c+d+\text{usw.}\) mit \(s\), so gilt \(K''\cos\frac{1}{2}s+L''\sin\frac{1}{2}s=0\), woraus folgt
\[
\begin{aligned}
h&=-\frac{K'\cos\frac{1}{2}s+L'\sin\frac{1}{2}s}
              {k\cos\frac{1}{2}s+l\sin\frac{1}{2}s}\\[5pt]
K''&=\phantom{-}\frac{K'l-L'k}
              {k\cos\frac{1}{2}s+l\sin\frac{1}{2}s}\sin\tfrac{1}{2}s\\[5pt]
L''&=-\frac{K'l-L'k}
              {k\cos\frac{1}{2}s+l\sin\frac{1}{2}s}\cos\tfrac{1}{2}s
\end{aligned}
\]
Ebenso leicht lässt sich zeigen, dass
\[
k=\pm\frac{\cos\frac{1}{2}s}{2^{\mu-1}},\quad
l=\pm\frac{\sin\frac{1}{2}s}{2^{\mu-1}}
\]
wobei das obere Zeichen gilt, wenn \(\frac{1}{2}\mu\) gerade ist, das untere, wenn es ungerade ist. Daher wird in den eben angegebenen Werten für \(K''\) und \(L''\) der Nenner \(=\pm\frac{1}{2^{\mu-1}}\), und somit gilt
\[
\begin{aligned}
K''&=\phantom{-}(K'\sin\tfrac{1}{2}s-L'\cos\tfrac{1}{2}s)\sin\tfrac{1}{2}s\\[5pt]
L''&=-(K'\sin\tfrac{1}{2}s-L'\cos\tfrac{1}{2}s)\cos\tfrac{1}{2}s
\end{aligned}
\]

\Needspace{46\baselineskip}
\InterpolationArticle{16}

Bisher haben wir Funktionen betrachtet, in denen sowohl Kosinus als auch Sinus vorkommen. Sehr häufig fehlen jedoch entweder diese oder jene; es ist zweckmäßig, diese Arten von Funktionen gesondert zu behandeln. Beide Fälle ließen sich zwar auf den bisher betrachteten allgemeinen Fall zurückführen; es scheint jedoch noch zweckmäßiger, diese Untersuchung auf das Problem von Artikel 3 zurückzuführen.

\textit{Erstens}, steht fest, dass alle Koeffizienten \(\beta',\beta'',\beta'''\) usw. in der Funktion \(X\) (Artikel 10) \(=0\) sind, so vermindert sich die Anzahl der Unbekannten auf \(m+1\), und es genügt daher, für ebenso viele verschiedene Werte von \(x\) die entsprechenden Werte der Funktion \(X\) zu kennen. Da sich in diesem Fall der Wert der Funktion \(X\) nicht ändert, wenn für \(x\) \(-x\) eingesetzt wird, dürfen offenbar nicht nur solche Werte von \(x\) als verschieden gelten, deren Differenz einem Umfang oder einem Vielfachen davon gleich ist, sondern auch solche, deren Summe \(0\) oder einem Umfang oder einem Vielfachen davon gleich ist. Allgemein müssen zwei Werte von \(x\), damit sie als verschieden gelten dürfen, un\SourceLeaf{300}%
gleiche Kosinuswerte haben. In diesem Sinn seien die \(m+1\) gegebenen Werte der Funktion \(X\) die Größen \(A, B, C, D\) usw., die den Werten von \(x\), nämlich \(a, b, c, d\) usw., entsprechen; gesucht ist ein allgemeiner Ausdruck für den Wert \(T\), den die Funktion für \(x=0\) annimmt.
Da sich die Funktion \(X\) in unserem Fall auch auf die Form
\[
\gamma+\gamma'\cos x+\gamma''(\cos x)^2+\gamma'''(\cos x)^3+\text{usw.}+\gamma^m(\cos x)^m
\]
bringen lässt, gilt nach Artikel 3
\[
\begin{aligned}
T={}&\frac{(\cos t-\cos b)(\cos t-\cos c)(\cos t-\cos d)\mathinner{.\,.\,.}}{(\cos a-\cos b)(\cos a-\cos c)(\cos a-\cos d)\mathinner{.\,.\,.}}\,A\\[5pt]
&+\frac{(\cos t-\cos a)(\cos t-\cos c)(\cos t-\cos d)\mathinner{.\,.\,.}}{(\cos b-\cos a)(\cos b-\cos c)(\cos b-\cos d)\mathinner{.\,.\,.}}\,B\\[5pt]
&+\frac{(\cos t-\cos a)(\cos t-\cos b)(\cos t-\cos d)\mathinner{.\,.\,.}}{(\cos c-\cos a)(\cos c-\cos b)(\cos c-\cos d)\mathinner{.\,.\,.}}\,C\\[5pt]
&+\frac{(\cos t-\cos a)(\cos t-\cos b)(\cos t-\cos c)\mathinner{.\,.\,.}}{(\cos d-\cos a)(\cos d-\cos b)(\cos d-\cos c)\mathinner{.\,.\,.}}\,D\\[5pt]
&+\text{usw.}
\end{aligned}
\]

\Needspace{46\baselineskip}
\textit{Zweitens}, verschwinden alle Koeffizienten \(\alpha,\alpha',\alpha'',\alpha'''\) usw., so beträgt die Anzahl der Unbekannten \(m\), und es genügt daher, die Werte der Funktion \(X\) für ebenso viele verschiedene Werte von \(x\) zu kennen. Da der Wert der Funktion \(X\) für \(x=-t\) aus dem Wert für
\(x=+t\) unmittelbar folgt (durch bloße Umkehrung des Vorzeichens), können diese Werte auch hier nicht als verschieden gelten. Daraus folgt leicht, dass hier wie im vorangehenden Fall nur solche Werte von \(x\) als verschieden gelten, deren Kosinuswerte ungleich sind. Da ferner der Wert der Funktion \(X\) für \(x=0,180^0,360^0\)
usw. und allgemein für einen Wert von \(x\), dessen Sinus \(=0\) ist, von selbst \(=0\) wird
und somit schon durch die Natur des Problems gegeben ist, setzen wir voraus, dass ein solcher Wert nicht unter den \(m\) gegebenen Werten mitgezählt wird. Da nun bekanntlich
\[
\frac{\sin nx}{\sin x}=2\cos(n-1)x+2\cos(n-3)x+2\cos(n-5)x+\mathinner{.\,.\,.}+2\cos x
\]
für einen geraden Wert von \(n\) und
\[
=2\cos(n-1)x+2\cos(n-3)x+2\cos(n-5)x+\mathinner{.\,.}+2\cos 2x+1
\]
für einen ungeraden Wert von \(n\) gilt, folgt leicht, dass \(\frac{X}{\sin x}\) auf die Form
\[
\delta+\delta'\cos x+\delta''(\cos x)^2+\delta'''(\cos x)^3+\text{usw.}+\delta^{m-1}(\cos x)^{m-1}
\]
\SourceLeaf{301}
gebracht werden kann. Da die Werte der Funktion \(\frac{X}{\sin x}\) für \(x=a,b,c,d\mathinner{.\,.\,.}\)
\(=\frac{A}{\sin a},\frac{B}{\sin b},\frac{C}{\sin c},\frac{D}{\sin d}\mathinner{.\,.\,.\,.}\) werden, ergibt sich daher nach demselben Artikel 3
\[
\begin{aligned}
T={}&\frac{\sin t(\cos t-\cos b)(\cos t-\cos c)(\cos t-\cos d)\mathinner{.\,.\,.}}{\sin a(\cos a-\cos b)(\cos a-\cos c)(\cos a-\cos d)\mathinner{.\,.\,.}}\,A\\[5pt]
&+\frac{\sin t(\cos t-\cos a)(\cos t-\cos c)(\cos t-\cos d)\mathinner{.\,.\,.}}{\sin b(\cos b-\cos a)(\cos b-\cos c)(\cos b-\cos d)\mathinner{.\,.\,.}}\,B\\[5pt]
&+\frac{\sin t(\cos t-\cos a)(\cos t-\cos b)(\cos t-\cos d)\mathinner{.\,.\,.}}{\sin c(\cos c-\cos a)(\cos c-\cos b)(\cos c-\cos d)\mathinner{.\,.\,.}}\,C\\[5pt]
&+\frac{\sin t(\cos t-\cos a)(\cos t-\cos b)(\cos t-\cos c)\mathinner{.\,.\,.}}{\sin d(\cos d-\cos a)(\cos d-\cos b)(\cos d-\cos c)\mathinner{.\,.\,.}}\,D\\[5pt]
&+\text{usw.}
\end{aligned}
\]
Im Übrigen ist in beiden Fällen die Formel für \(T\) mit der Funktion \(X\) identisch, wenn
\(t\) in \(x\) geändert wird, da jene ohne Einschränkung für jeden beliebigen Wert von \(t\) gilt.

\endgroup
